\chapter{Continuous Chance}

\begin{center}\itshape
A uniform point of $[0,1]$ equals $\frac12$ with probability $0$, and it equals every other\\
single number with probability $0$. The average is an integral, not a sum.
\end{center}

Chapters~1 to~4 added probabilities of intervals and of finite-bit events. A single point inside $[0,1)$ has length $0$. If a measurement can fall anywhere on a line, its expectation cannot be a sum over values. This chapter builds that expectation from the Riemann integral of Volume~1, proves the repeated integral of a continuous function on a rectangle, and constructs as many independent normal variables as a later proof asks for. It does not assign a length to every subset of the line. Vitali showed in 1905 that countable additivity, translation invariance, and a length for every subset cannot hold together, if one grants the axiom of choice. The sets used below are intervals and finite unions of intervals, together with certain limits of such events whose probabilities are established by explicit approximations. We do not assume a general length or probability for every zero set of a continuous function.

\section{Distribution functions and densities}

Let $X$ be a random variable, in the sense that $\{X\le x\}$ is an event whose probability we can compute, for every real $x$. The \emph{distribution function} is $F(x)=P(X\le x)$. It is increasing, $F(x)\to0$ as $x\to-\infty$ and $F(x)\to1$ as $x\to\infty$ whenever the model is a probability, and it is right-continuous when the events $\{X\le x+1/n\}$ shrink to $\{X\le x\}$ with probabilities tending to $F(x)$. For the variables in this chapter that limit is part of the construction.

\begin{definition}[density]\label{def:density}
A nonnegative function $f$ that is continuous except at finitely many points (and has finite one-sided limits there) is a \emph{density} of $X$ when
\[
 P(a<X\le b)=\int_a^b f(t)\,\dd t
\]
for all $a<b$, the integral being the improper Riemann integral of Volume~1 if the limits are infinite, and $\int_{-\infty}^{\infty}f=1$.
\end{definition}

If two functions in this class both satisfy the identity, they agree at every point where both are continuous. A change on a finite set does not alter the Riemann integral, so the identity does not pin down the values at the finitely many jumps. At a continuity point the fundamental theorem says that the derivative of $\int_0^x f$ is $f(x)$.

The uniform law on $[0,1]$ has density $1$ on that interval and $0$ outside. An exponential random variable of rate $\lambda>0$ has density $\lambda e^{-\lambda x}$ for $x>0$. Its integral is $1$, and $\mathbb E[X]=\int_0^{\infty}x\lambda e^{-\lambda x}\,\dd x=1/\lambda$, by parts. The \emph{standard normal} density is
\[
 \phi(x)=\frac{1}{\sqrt{2\pi}}e^{-x^2/2}.
\]
Volume~1, Chapter~8, gives $\int_0^{\infty}e^{-u^2}\,\dd u=\sqrt{\pi}/2$. The integrand is even, so the integral over the whole line is $\sqrt{\pi}$. The substitution $u=x/\sqrt{2}$ turns that into $\int_{-\infty}^{\infty}e^{-x^2/2}\,\dd x=\sqrt{2\pi}$, so $\phi$ integrates to $1$.

\begin{proposition}[expectation]\label{prop:densitymean}
If $X$ has density $f$ and $\int |x|f(x)\,\dd x<\infty$, then $\mathbb E[X]=\int_{-\infty}^{\infty} x f(x)\,\dd x$. If $g$ is continuous and $\int |g|f<\infty$, then $\mathbb E[g(X)]=\int g(x)f(x)\,\dd x$.
\end{proposition}
\begin{proof}
For a variable with a density, define $\mathbb E[g(X)]$ to be $\int g f$ whenever $g$ is continuous and $\int|g|f<\infty$. The definition has to agree with a finite sum when $g$ is a step.

Let $g$ be constant on finitely many intervals and zero outside them, and put every discontinuity of $f$ among the partition points. On each piece the probability is the integral of $f$, so the sum of value times probability equals $\int g f$.

Now let $g$ be continuous on a compact interval $[A,B]$, zero outside, and let $f$ be bounded there. Uniform continuity gives a step function $s$, constant between the jumps of $f$, with $|g-s|<\varepsilon$ on $[A,B]$. Then
\[
 \Bigl|\int g f-\int s f\Bigr|\le\varepsilon\int_A^B f\le\varepsilon,
\]
and the step integral is a finite sum of values times probabilities. So the integral of $g f$ is a limit of expectations of explicit finite experiments.

For a general continuous $g$ with $\int|g|f<\infty$, choose $A$ so that $\int_{|x|>A}|g|f<\varepsilon$. On $[-A,A]$ approximate $g$ by a step function $s$ that is zero outside $[-A,A]$, with the finitely many jumps of $f$ used as partition points. Then $\int s f$ is a finite sum of values times probabilities, and it lies within $\varepsilon$ of the integral of $g f$ over $[-A,A]$. The integral of $|g|f$ over the two tails is less than $\varepsilon$. The integral over the line is therefore within $2\varepsilon$ of that finite sum. There is no earlier sum that defines $\mathbb E[g(X)]$ for a continuous $g$. The integral is the definition, and the step comparison is the check that it extends the finite sum.
\end{proof}

\begin{proposition}[Gaussian tail]\label{prop:gausstail}
For $x>0$,
\[
 \int_x^{\infty}e^{-t^2/2}\,\dd t\le\frac1x e^{-x^2/2}.
\]
\end{proposition}
\begin{proof}
On $[x,\infty)$ one has $1\le t/x$, so the integrand is at most $(t/x)e^{-t^2/2}$. The antiderivative of $t e^{-t^2/2}$ is $-e^{-t^2/2}$.
\end{proof}

The same integration gives the mean and the variance. Let $Z$ have density $\phi$. The integrand $x\phi(x)$ is odd. It is absolutely integrable, because
\[
 \int_0^{\infty} x e^{-x^2/2}\,\dd x=\bigl[-e^{-x^2/2}\bigr]_0^{\infty}=1,
\]
so $\mathbb E|Z|=2/\sqrt{2\pi}=\sqrt{2/\pi}$ and the signed integral is $0$. Thus $\mathbb E[Z]=0$. For the second moment integrate by parts on $[0,A]$, with $u=x$ and $\mathrm dv=x e^{-x^2/2}\,\dd x$:
\[
 \int_0^A x^2 e^{-x^2/2}\,\dd x=\bigl[-x e^{-x^2/2}\bigr]_0^A+\int_0^A e^{-x^2/2}\,\dd x.
\]
The boundary term is $A e^{-A^2/2}$. It tends to $0$: the series for $e^y$ gives $e^{A^2/4}\ge (A^2/4)^2/2=A^4/32$ for every $A>0$, so $A e^{-A^2/2}=A e^{-A^2/4}/e^{A^2/4}\le 32/A^3$. The remaining integral tends to $\sqrt{2\pi}/2$. Doubling and dividing by $\sqrt{2\pi}$ leaves
\[
 \mathbb E[Z^2]=\frac{2}{\sqrt{2\pi}}\cdot\frac{\sqrt{2\pi}}{2}=1.
\]
Since the mean is $0$, the variance is $1$. Chapter~6 uses this value when it scales a sum to the density $\phi$.

\section{Changing the variable}

A density is a recipe for interval probabilities. A smooth function of the variable pushes those intervals forward, and the fundamental theorem brings the density back.

\begin{examplex}[an exponential from one uniform]
Let $U$ be uniform on $(0,1)$ and set $X=-\ln U$. For $x>0$,
\[
 P(X>x)=P(U<e^{-x})=e^{-x},
\]
because $U$ has no atom at the endpoint. The distribution function on $(0,\infty)$ is $F(x)=1-e^{-x}$. Differentiating by the fundamental theorem gives the density $e^{-x}$. Integration by parts, Volume~1, Chapter~8, gives the mean:
\[
 \mathbb E[X]=\int_0^{\infty} x e^{-x}\,\dd x=\bigl[-xe^{-x}\bigr]_0^{\infty}+\int_0^{\infty}e^{-x}\,\dd x=1.
\]
The boundary term at infinity vanishes because $x e^{-x}\to0$, which is the standard limit from Volume~1. Scaling gives the rate. If $Y=X/\lambda$ for $\lambda>0$, then $P(Y>y)=e^{-\lambda y}$ and the density is $\lambda e^{-\lambda y}$, with mean $1/\lambda$. Chapter~10 starts from this tail.
\end{examplex}

\begin{examplex}[the mean as an integral of the tail]
Suppose $X\ge0$ has a density $f$ and $\int_0^{\infty} x f(x)\,\dd x<\infty$. Then
\[
 \mathbb E[X]=\int_0^{\infty} P(X>t)\,\dd t.
\]
Let $F(t)=\int_0^t f$. For $0\le t\le A$,
\[
 P(X>t)=\int_t^A f+\int_A^{\infty} f=F(A)-F(t)+P(X>A).
\]
Integrate in $t$ from $0$ to $A$. The constant terms contribute $A P(X>A)$. Integration by parts gives
\[
 \int_0^A\bigl(F(A)-F(t)\bigr)\,\dd t=A F(A)-\Bigl[t F(t)\Bigr]_0^A+\int_0^A t f(t)\,\dd t=\int_0^A t f(t)\,\dd t.
\]
The remainder $A P(X>A)$ is at most $\int_A^{\infty} x f(x)\,\dd x$, which tends to $0$ as $A\to\infty$ because the mean integral converges. Therefore the tail integral equals the mean.

Check it on the exponential of rate $1$: $\int_0^{\infty} e^{-t}\,\dd t=1$, which is the mean just computed.
\end{examplex}

\section{Repeated integrals}

Volume~1 integrates in one variable. A joint density needs two.

\begin{proposition}[continuous Fubini on a rectangle]\label{prop:fubini}
Let $f$ be continuous on the rectangle $[a,b]\times[c,d]$. Then
\[
 \int_a^b\Bigl(\int_c^d f(x,y)\,\dd y\Bigr)\dd x
 =\int_c^d\Bigl(\int_a^b f(x,y)\,\dd x\Bigr)\dd y.
\]
Both inner integrals are continuous in the outer variable.
\end{proposition}
\begin{proof}
Continuity on a compact rectangle is uniform. The Riemann sums over a grid are finite double sums, and addition is commutative, so the two orders of summation agree exactly at each grid. As the mesh tends to $0$, both iterated integrals equal the common limit of those sums. The inner integral is continuous in the outer variable because $|f|$ is bounded and the integrand moves uniformly when the outer variable moves a little.
\end{proof}

This is the rectangle version of Fubini used in this volume. It is not the theorem for arbitrary nonnegative measurable functions, which Volume~7 proves. For an improper integral we use it on compact rectangles first and pass to a limit only when the tails are controlled, for example by the explicit Gaussian decay proved above. The existence of two iterated improper integrals alone would not justify exchanging their order.

\begin{proposition}[polar rectangles]\label{prop:polar}
Let $h$ be continuous and bounded on $\R^2$. For a fixed radius $a>0$,
\[
 \int\!\!\int_{x^2+y^2\le a^2} h(x,y)\,\dd x\,\dd y
 =\int_0^{2\pi}\int_0^a h(r\cos\theta,r\sin\theta)\,r\,\dd r\,\dd\theta,
\]
where the left side means the iterated integral of $h$ over the disk, obtained as the limit of integrals of continuous functions that agree with $h$ on the closed disk of radius $a-\eta$ and vanish outside the disk of radius $a+\eta$. Proposition~\ref{prop:fubini} applies to each of those continuous integrands on a large square. The symmetric difference of the two disks has area at most $2\pi a\eta+\pi\eta^2$, so the boundedness of $h$ makes the limit unambiguous. The error in replacing that integral by a union of polar rectangles tends to $0$ with the mesh. The letter $a$ is not a random radius.
\end{proposition}
\begin{proof}
Cut $0\le r\le a$ and $0\le\theta\le 2\pi$ into small pieces of lengths $\Delta r$ and $\Delta\theta$. The region $r_i\le r\le r_i+\Delta r$, $\theta_j\le\theta\le\theta_j+\Delta\theta$ is a polar rectangle. Its area is
\[
 \tfrac12\bigl((r_i+\Delta r)^2-r_i^2\bigr)\Delta\theta=(r_i\Delta r+\tfrac12(\Delta r)^2)\Delta\theta.
\]
Summing the second terms over a partition gives at most $\pi a\cdot\Delta r$, which tends to $0$ with the mesh. On each polar rectangle, $h$ varies by less than a prescribed $\varepsilon$ once the mesh is small, by uniform continuity on the disk. The Riemann sum with weights $r_i\Delta r\Delta\theta$ therefore differs from the Cartesian integral of $h$ by at most $\varepsilon$ times the area $\pi a^2$, plus a constant times the area of the symmetric difference, plus the $(\Delta r)^2$ contribution. All three can be made small.
\end{proof}

The factor $r$ is the area of a thin ring, divided by its thickness and by the angle. It is not an extra assumption.

\begin{figure}[ht]
\centering
\fig{ch05-polar}
\caption{One polar cell, from radius $r$ to $r+\Delta r$ and through an angle $\Delta\theta$. The outer arc has length $(r+\Delta r)\Delta\theta$. The area of the cell is $(r\Delta r+\tfrac12(\Delta r)^2)\Delta\theta$, which is the weight in the proof. The shaded cell is a schematic of that region.}
\label{fig:polar}
\end{figure}

\section{Independent copies from digits}

Chapter~1 built one uniform $U=\sum_{m\ge1}b_m 2^{-m}$ from fair bits. Countably many independent uniforms need a partition of the bits.

\begin{proposition}[independent uniforms]\label{prop:manyuniforms}
Enumerate the bit positions as $2^{k-1}(2m-1)$ for integers $k\ge1$ and $m\ge1$. Let $U_k$ be the number whose binary digits are the bits in the $k$th block, in order of $m$. On the exceptional sequence where all digits of a block are $1$, set that $U_k=0$ rather than $1$; the exception has outer content $0$, because it lies in the cylinder of the first $N$ ones for every $N$. Then each $U_k$ takes values in $[0,1)$ and is uniform on that interval, and every finite subcollection is independent: the joint probability of finitely many interval conditions is the product of the lengths.
\end{proposition}
\begin{proof}
Each block is an infinite set of distinct bit positions, and distinct blocks are disjoint. Any condition on the first $N$ digits of $U_{k_1},\ldots,U_{k_r}$ is a condition on finitely many bits. Proposition~\ref{prop:bits} says those bits are fair and independent, so the probability is $2^{-rN}$. That is the volume of a product of dyadic intervals. For a product of general intervals, choose dyadic intervals inside each factor and dyadic intervals containing it, with total length gaps less than any prescribed $\varepsilon$. The cylinder probability is then trapped between two products of lengths that both approach the product of the target lengths. The joint probability is the product.
\end{proof}

\begin{examplex}[the smaller and the larger of two uniforms]
Let $U$ and $V$ be the first two of these independent uniforms, and write $M=\max(U,V)$ and $L=\min(U,V)$. For $t\in[0,1]$,
\[
 P(M\le t)=P(U\le t)P(V\le t)=t^2,
\]
because the event is a product of intervals. Differentiating gives the density $2t$ of the larger value. Likewise $P(L>s)=(1-s)^2$ for $s\in[0,1]$, so the smaller value has density $2(1-s)$. The means are
\[
 \mathbb E[M]=\int_0^1 2t\cdot t\,\dd t=\frac23,\qquad
 \mathbb E[L]=\int_0^1 2(1-s)s\,\dd s=\frac13.
\]
Exercise~\ref{ex:order} asks for the marginal density of the $k$th among $n$ uniforms. The joint law of the first two is the change of variables itself.
\end{examplex}

\begin{proposition}[the ordered pair]\label{prop:orderpair}
Let $U$ and $V$ be independent and uniform on $[0,1]$, and set $L=\min(U,V)$, $M=\max(U,V)$. The joint density of $(L,M)$ is $2$ on the triangle $0<s<t<1$ and $0$ elsewhere. The gap $M-L$ has density $2(1-u)$ on $(0,1)$ and mean $1/3$.
\end{proposition}
\begin{proof}
Take $0\le s<s+\Delta s\le t<t+\Delta t\le 1$, so the rectangle in the $(s,t)$ plane does not meet the diagonal. It has two preimages in the $(u,v)$ square: $u$ in the $s$-slot and $v$ in the $t$-slot, or the swap. Each preimage is a rectangle of area $\Delta s\,\Delta t$, and the two are disjoint. Independence gives each the probability $\Delta s\,\Delta t$. The joint probability is therefore $2\Delta s\,\Delta t$, which is the density $2$.

The gap exceeds $u\in(0,1)$ when $0<s<t<1$ and $t-s>u$. Integrate the density:
\[
 P(M-L>u)=\int_0^{1-u} 2\bigl(1-(s+u)\bigr)\,\dd s=(1-u)^2.
\]
Differentiate in $u$. The density of the gap is $2(1-u)$. The tail formula already checked for a nonnegative variable gives
\[
 \mathbb E[M-L]=\int_0^1(1-u)^2\,\dd u=\frac13,
\]
which is also $\mathbb E[M]-\mathbb E[L]=\tfrac23-\tfrac13$.
\end{proof}

\begin{proposition}[inverse distribution function]\label{prop:inverse}
Let $F$ be continuous and strictly increasing, with limits $0$ and $1$ at $\pm\infty$, and set $F^{-1}(u)=\inf\{x:F(x)\ge u\}$ for $u\in(0,1)$. If $U$ is uniform on $[0,1)$, define $F^{-1}(U)$ arbitrarily, say as $0$, on the event $\{U=0\}$ of probability $0$. With that harmless endpoint convention it is a random variable on the whole sample space and $P(F^{-1}(U)\le x)=F(x)$.
\end{proposition}
\begin{proof}
For $u\in(0,1)$, strict increase and continuity give $F(F^{-1}(u))=u$ and $F^{-1}(F(x))=x$. The inequality $F^{-1}(u)\le x$ is equivalent to $u\le F(x)$. A uniform number falls in $(0,F(x)]$ with probability $F(x)$, and the separately defined event $\{U=0\}$ has probability zero.
\end{proof}

Apply this to the standard normal distribution function $\Phi(x)=\int_{-\infty}^x\phi$. It is continuous and strictly increasing from $0$ to $1$. The variables $Z_k=\Phi^{-1}(U_k)$ are independent standard normals: every finite subcollection has the product distribution function, because the $U_k$ are independent and each coordinate transforms by $\Phi^{-1}$.

Theorem~\ref{thm:sllnbound} does \emph{not} apply directly to $X_k=F^{-1}(U_k)$: one such variable generally depends on infinitely many digits, whereas the theorem assumes a finite-digit function. Here is the additional step for a bounded nondecreasing transformation $Q$ of a uniform point, with finite endpoint values $Q(0)$ and $Q(1)$. (For a bounded law with a right-continuous distribution function $F$, the generalised inverse $Q(u)=\inf\{x:F(x)\ge u\}$ is such a transformation, after defining its endpoints by one-sided limits. Right-continuity gives $Q(u)\le x$ exactly when $u\le F(x)$, so it has law $F$.) Partition $[0,1)$ into intervals $I_{j,m}=[j2^{-m},(j+1)2^{-m})$ and, on each interval, define the step bounds $L_m(u)=Q(j2^{-m})$ and $R_m(u)=Q((j+1)2^{-m})$. Then $L_m(U_k)\le Q(U_k)\le R_m(U_k)$, and both bounds depend only on the first $m$ digits of $U_k$. Moreover their expected gap is exactly
\[
\mathbb E[R_m(U_k)-L_m(U_k)]=2^{-m}\sum_{j=0}^{2^m-1}\bigl(Q((j+1)2^{-m})-Q(j2^{-m})\bigr)=2^{-m}(Q(1)-Q(0)).
\]
For each fixed $m$, Chapter~4's strong law applies to the independent copies of $L_m(U_k)$ and $R_m(U_k)$. Intersecting the probability-one conclusions over all integers $m$ and squeezing the averages of $Q(U_k)$ between them gives convergence to its integral mean, since the expectation of the two bounds differs by at most $2^{-m}(Q(1)-Q(0))$. Thus the finite-digit hypothesis can be removed for this bounded quantile model by an explicit approximation, without borrowing a general product-measure theorem.

\section{Passing a limit under the integral}

\begin{theorem}[dominated convergence for continuous integrands]\label{thm:dct}
Let $f_n$ and $f$ be continuous on $\R$, with $f_n\to f$ uniformly on every bounded interval, and suppose $|f_n(x)|\le g(x)$ for a continuous $g$ with $\int_{-\infty}^{\infty}g<\infty$. Then $\int f_n\to\int f$.
\end{theorem}
\begin{proof}
The bound $|f|\le g$ passes to the limit, so the improper integral of $f$ exists absolutely. Fix $\varepsilon>0$ and choose $A$ so that $\int_{|x|>A}g<\varepsilon$. On $[-A,A]$ the convergence is uniform, so for large $n$ the integral of $|f_n-f|$ over that segment is less than $\varepsilon$. The two tails together contribute at most $2\varepsilon$, giving $|\int f_n-\int f|<3\varepsilon$. Since $\varepsilon$ is arbitrary, the difference tends to zero.
\end{proof}

Differentiation under the integral sign in Chapter~6 requires more than this theorem by itself: the relevant difference quotients must converge locally uniformly and admit one integrable continuous bound. When those hypotheses are verified, the theorem applies to the quotients and justifies taking their integrals to the derivative.

\section{Box--Muller}

\begin{theorem}[Box--Muller]\label{thm:boxmuller}
Let $U$ and $V$ be independent uniforms on $(0,1)$. Set $R=\sqrt{-2\ln U}$ and $\Theta=2\pi V$, and set $X=R\cos\Theta$, $Y=R\sin\Theta$. Then $X$ and $Y$ are independent standard normals.
\end{theorem}
\begin{proof}
The distribution function of $R$ is $P(\sqrt{-2\ln U}\le r)=P(U\ge e^{-r^2/2})=1-e^{-r^2/2}$ for $r>0$, so $R$ has density $r e^{-r^2/2}$ on $(0,\infty)$. The angle $\Theta$ is independent of $U$, hence of $R$, and uniform on $(0,2\pi)$. For a bounded continuous function $q$ on $\R^2$, Proposition~\ref{prop:polar} and the product density of $(R,\Theta)$ give
\begin{align*}
\mathbb E[q(X,Y)]
&=\int_0^{2\pi}\int_0^{\infty} q(r\cos\theta,r\sin\theta)\,\frac{1}{2\pi}e^{-r^2/2}\,r\,\dd r\,\dd\theta\\
&=\int\!\!\int_{\R^2} q(x,y)\,\frac{1}{2\pi}e^{-(x^2+y^2)/2}\,\dd x\,\dd y.
\end{align*}
The identity is first the polar formula on the disk of radius $T$, then the limit $T\to\infty$. Here $T$ is a cutoff, not the random radius $R$. The integral of the product density outside that disk is at most the Gaussian tail $\int_{r>T} r e^{-r^2/2}\,\dd r=e^{-T^2/2}$, which tends to $0$. The Jacobian factor $r$ cancels the $r$ in the density of $R$. The remaining function of $(x,y)$ is the product $\phi(x)\phi(y)$. Therefore the joint density of $(X,Y)$ is the product of two standard normal densities.
\end{proof}

One uniform would have been enough, by Proposition~\ref{prop:inverse}. Box--Muller makes the independence visible without inverting $\Phi$.

The radius $R$ has tail $P(R>r)=e^{-r^2/2}$ for $r>0$, from the distribution function in the proof. The tail formula above therefore gives
\[
 \mathbb E[R]=\int_0^{\infty}e^{-r^2/2}\,\dd r=\sqrt{\frac{\pi}{2}},
\]
the half-line integral of Volume~1, Chapter~8, after the substitution $u=r/\sqrt{2}$ in $\int_0^{\infty}e^{-u^2}\,\dd u=\sqrt{\pi}/2$.

\begin{figure}[ht]
\centering
\fig{ch05-cantor}
\caption{The first two stages of the middle-thirds construction. Each stage removes the open middle third of every remaining interval. The total length left after $n$ stages is $(2/3)^n$, so the intersection can be covered by intervals of arbitrarily small total length. Exercise~\ref{ex:cantor} puts a random series on what remains.}
\label{fig:cantor}
\end{figure}

Sums of independent copies of a normal are again normal, by Lemma~\ref{lem:normalsum} in the next chapter. The shape of a general sum is that chapter's theorem.

% Continuous sampling models, input from Chapter 5.
\section{Bertrand's three experiments}

A chord in a circle is longer than the side of an inscribed equilateral triangle when its midpoint lies inside the concentric disk of half the radius. Three different experiments give three probabilities. Joseph Bertrand set the question in \emph{Calcul des probabilit\'es} (1889), page~4, and refused to choose. Take the radius to be $1$. The side of the inscribed equilateral triangle subtends the angle $2\pi/3$ at the centre.

\begin{figure}[ht]
\centering
\fig{ch01-bertrand}
\caption{Three experiments, one circle of radius $1$. Left: one endpoint is fixed and the shorter angle $\alpha$ to the other endpoint exceeds $2\pi/3$ on a third of the circle. Middle: the distance $D$ from the centre to the chord is drawn uniformly from $[0,1]$, and the long chords are those with $D<1/2$. Right: the midpoint is uniform in the disk, and the long chords are those whose midpoint falls in the inner disk, of a quarter of the area.}
\label{fig:bertrand}
\end{figure}

\begin{theorem}[three models for a chord]\label{thm:bertrand}
\begin{enumerate}[label=(\alph*)]
\item If the two endpoints are independent and uniform on the circle, the probability that the chord is long is $\frac13$.
\item If a radius is chosen, and the distance from the centre to the chord is then uniform on $[0,1]$ and perpendicular to that radius, the probability is $\frac12$.
\item If the midpoint is uniform in the disk, the probability is $\frac14$.
\end{enumerate}
\end{theorem}
\begin{proof}
(a) Rotational symmetry lets us fix one endpoint. The other endpoint is uniform on the circle. Let $\alpha\in(0,\pi]$ be the shorter angle between the endpoints. For $\theta\in(0,\pi]$, the event $\alpha\le\theta$ means that the free endpoint lies in one of the two arcs of length $\theta$ adjacent to the fixed endpoint, a fraction $\theta/\pi$ of the circle. So $P(\alpha\le\theta)=\theta/\pi$, and $\alpha$ has density $1/\pi$ on $(0,\pi]$. The chord is long when $\alpha>2\pi/3$, which has probability $(\pi-2\pi/3)/\pi=1/3$.

(b) In this model the distance $D$ is uniform on $[0,1]$ by the description of the experiment, not by a theorem about endpoints. The chord is long when $D<1/2$, which has probability $1/2$.

(c) The set of midpoints that give a long chord is the open disk of radius $1/2$. Its area is $\pi/4$. The whole disk has area $\pi$. A uniform midpoint falls in the inner disk with probability $1/4$. The area of a disk is the integral computed in Volume~1, Chapter~8: the quarter circle $y=\sqrt{1-x^2}$ on $[0,1]$ has area $\pi/4$.
\end{proof}

The three numbers disagree because the three sets of outcomes disagree. Model (a) spreads endpoints. Model (b) spreads distances along a radius. Model (c) spreads area. Bertrand's sentence is the instruction this book follows. Write the experiment, then compute.

\begin{examplex}[the length, under one model]
Keep the radius equal to $1$, and use model (b): the distance $D$ from the centre to the chord is uniform on $[0,1]$. The chord length is $L=2\sqrt{1-D^2}$. It falls from $2$ to $0$ as $D$ runs from $0$ to $1$. For $\ell\in(0,2)$,
\[
 P(L\le\ell)=P\Bigl(D\ge\sqrt{1-(\ell/2)^2}\Bigr)=1-\sqrt{1-(\ell/2)^2}.
\]
Differentiate. The chain rule gives the density
\[
 f_L(\ell)=\frac{\ell}{4\sqrt{1-(\ell/2)^2}}\qquad\text{on }(0,2).
\]
The long chords are those longer than the side $\sqrt3$ of the inscribed equilateral triangle. The formula yields $P(L>\sqrt3)=1/2$, which is clause (b) of the theorem again: $\sqrt{1-3/4}=1/2$.

The density blows up as $\ell\uparrow2$ and tends to $0$ as $\ell\downarrow0$. Uniform distance spends as much probability on $D<1/2$ as on $D>1/2$, and $D<1/2$ is exactly the long chords. Area measure does not. The disk of radius $1/2$ has a quarter of the area, which is why model (c) gives probability $1/4$ for the same geometric event. The map from distance to length is not the identity, so a uniform law for $D$ is not a uniform law for $L$, and neither of those is the area law for the midpoint.
\end{examplex}

\section{Buffon's needle}

Parallel lines are ruled on the plane at distance $d>0$. A needle of length $\ell$, with $0<\ell\le d$, is thrown. Buffon described this experiment in the \emph{Essai d'arithm\'etique morale}, in the supplement to the \emph{Histoire naturelle}, in 1777. The model we adopt is the following. Let $X$ be the distance from the needle's centre to the nearest line, uniform on $[0,d/2]$. Let $\Theta$ be the angle between the needle and the direction of the lines, taken in $[0,\pi]$ by turning the needle end for end, and take $\Theta$ uniform on $[0,\pi]$ and independent of $X$. The needle meets a line when the vertical half-projection reaches the nearest line, that is when $X\le(\ell/2)\sin\Theta$. The restriction $\ell\le d$ keeps the needle from meeting two lines at once, so this inequality is the whole event.

\begin{figure}[ht]
\centering
\fig{ch01-buffon}
\caption{Two rulings at distance $d$. The needle has length $\ell\le d$. It meets a ruling when the distance $X$ from its centre to the nearest ruling is at most $(\ell/2)\sin\Theta$.}
\label{fig:buffon}
\end{figure}

\begin{theorem}[Buffon's needle]\label{thm:buffon}
Under the model just stated, the probability that the needle meets a line is $2\ell/(\pi d)$.
\end{theorem}
\begin{proof}
Independence and uniformity put density $2/d$ on $X\in[0,d/2]$ and density $1/\pi$ on $\Theta\in[0,\pi]$. The probability is the integral of the product density over the region $0\le x\le(\ell/2)\sin\theta$. For each fixed $\theta$ the inner integral in $x$ is an ordinary one-variable integral, and $(\ell/2)\sin\theta\le d/2$ because $\ell\le d$ and $\sin\theta\le1$. So
\begin{align*}
P
&=\int_0^\pi\frac1\pi\Biggl(\int_0^{(\ell/2)\sin\theta}\frac2d\,\mathrm dx\Biggr)\mathrm d\theta
=\frac{2}{d\pi}\int_0^\pi\frac\ell2\sin\theta\,\mathrm d\theta
=\frac{\ell}{d\pi}\bigl[-\cos\theta\bigr]_0^\pi
=\frac{2\ell}{\pi d}.
\end{align*}
The iterated integral of a continuous function on a rectangle, cut down to a region described by $\sin$, is justified by writing it as a single integral of a continuous function of $\theta$, which Volume~1, Chapter~8, allows. The repeated-integral statement for a continuous integrand on a rectangle is Proposition~\ref{prop:fubini} above.
\end{proof}

When $\ell=d$, the probability is $2/\pi\approx0.6366$. A needle as long as the gap hits a line more often than not, and not twice as often as a needle of half the length would suggest if one forgot the angle. The factor $\sin\Theta$ is the whole geometry.

The same integral is the first geometric probability in the series. Volume~9 returns to it for a curve rather than a needle. The proof there may use more geometry. It does not replace the computation above.


\section{Exercises}

\begin{exercise}[$\star$]\label{ex:counter}
For $n\ge2$ let $f_n$ be the tent of height $n$ and width $2/n$ centred at $0$: $f_n(0)=n$, $f_n$ is linear from $(-1/n,0)$ to $(0,n)$ and from $(0,n)$ to $(1/n,0)$, and $f_n=0$ outside. Show that $\int f_n=1$ for every $n$, while $f_n(x)\to0$ for every $x\ne0$. Why does Theorem~\ref{thm:dct} not apply?
\end{exercise}

\begin{exercise}[$\star$]\label{ex:densityunique}
Suppose $f$ and $g$ are continuous and $\int_a^b f=\int_a^b g$ for every $a<b$. Show that $f=g$.
\end{exercise}

\begin{exercise}[$\star\star$]\label{ex:tail}
Prove Proposition~\ref{prop:gausstail}, and deduce $P(Z>x)\le e^{-x^2/2}/(x\sqrt{2\pi})$ for a standard normal $Z$ and $x>0$.
\end{exercise}

\begin{exercise}[$\star\star$]\label{ex:cauchy}
Show that $f(x)=1/(\pi(1+x^2))$ is a density. Show that $\int |x|f(x)\,\dd x=\infty$, so the recipe in Proposition~\ref{prop:densitymean} does not define a mean.
\end{exercise}

\begin{exercise}[$\star\star$]\label{ex:gamma}
For each integer $n\ge1$, set $I_n=\int_0^{\infty}x^{n-1}e^{-x}\,\dd x$. Show that $I_1=1$ and $I_{n+1}=n I_n$, hence $I_n=(n-1)!$.
\end{exercise}

\begin{exercise}[$\star\star$]\label{ex:order}
Let $U_1,\ldots,U_n$ be independent uniforms on $[0,1]$, and let $U_{(k)}$ be the $k$th smallest. Show that $U_{(k)}$ has density
\[
 n\binom{n-1}{k-1}t^{k-1}(1-t)^{n-k}\qquad\text{on }[0,1],
\]
and that $\mathbb E[U_{(k)}]=k/(n+1)$. You may use Exercise~\ref{ex:gamma} to evaluate the beta integral.
\end{exercise}

\begin{exercise}[$\star\star$]\label{ex:cantor}
Let $B_1,B_2,\ldots$ be independent fair bits and $X=\sum_{k\ge1}2B_k 3^{-k}$. Let $C$ be the set of numbers in $[0,1]$ that have a ternary expansion using only the digits $0$ and $2$.
\begin{enumerate}[label=(\alph*)]
\item Show that the total length of the intervals that remain at stage $n$ of the middle-thirds construction is $(2/3)^n$, so $C$ can be covered by intervals of arbitrarily small total length.
\item Show that $P(X\in C)=1$ and that $P(X=x)=0$ for every single $x$.
\item Show that the distribution function of $X$ is continuous, and that $X$ has no continuous density.
\end{enumerate}
\end{exercise}

\begin{exercise}[$\star\star$]\label{ex:usquare}
Let $U$ be uniform on $(0,1)$ and set $Y=U^2$. Find the distribution function and the density of $Y$ on $(0,1)$, and check that $\mathbb E[Y]=\int_0^1 u^2\,\dd u$.
\end{exercise}

\begin{exercise}[$\star$]\label{ex:radius}
In the Box--Muller construction, show that $R^2=X^2+Y^2$ has density $\tfrac12 e^{-r/2}$ on $(0,\infty)$, the exponential law of rate $1/2$.
\end{exercise}

\section{Hints}

\begin{hint}{ex:counter}
The area of a triangle is half base times height. Any dominant that lies above every tent is at least of height $1/|x|$ near $0$.
\end{hint}
\begin{hint}{ex:densityunique}
Differentiate $\int_0^x(f-g)$ by the fundamental theorem.
\end{hint}
\begin{hint}{ex:tail}
Compare the integrand with $(t/x)e^{-t^2/2}$.
\end{hint}
\begin{hint}{ex:cauchy}
The arctangent is an antiderivative. Compare $|x|/(1+x^2)$ with $1/(2|x|)$ for large $|x|$.
\end{hint}
\begin{hint}{ex:gamma}
Integrate by parts with $u=x^n$ and $\mathrm dv=e^{-x}\,\mathrm dx$.
\end{hint}
\begin{hint}{ex:order}
The probability that exactly $k-1$ samples fall below $t$, one falls in a small interval at $t$, and the rest fall above, is a multinomial count. For the mean, write the integral as a beta value and use Exercise~\ref{ex:gamma}.
\end{hint}
\begin{hint}{ex:cantor}
The ternary expansion of $X$ is $0.(2B_1)(2B_2)\ldots$ in base $3$. A continuous density would vanish on every removed middle third, and those intervals meet every subinterval of $[0,1]$.
\end{hint}
\begin{hint}{ex:usquare}
$P(Y\le y)=P(U\le\sqrt y)$ for $y\in(0,1)$. Differentiate, then integrate $y$ times the density.
\end{hint}
\begin{hint}{ex:radius}
$U=e^{-R^2/2}$, and $U$ is uniform.
\end{hint}

\section{Solutions}

\begin{solution}{ex:counter}
\idea A tall thin tent has area $1$ and disappears at each fixed $x\ne0$. No single integrable function sits above all of them.
\writeup The graph consists of two right triangles of base $1/n$ and height $n$, or one triangle of base $2/n$ and height $n$. Its area is $\tfrac12\cdot(2/n)\cdot n=1$. For $x\ne0$, eventually $1/n<|x|$, so $f_n(x)=0$. At $0$, $f_n(0)=n$ does not converge to $0$, but a single point does not affect the integral of the limit if one sets the limit equal to $0$ everywhere. If a continuous $g$ dominated every tent, it would satisfy $g(0)\ge f_n(0)=n$ for every $n$, which is impossible for a real-valued continuous function. Moreover, $f_n$ does not converge to a finite value at $0$ and does not converge uniformly on any interval containing $0$. Thus neither the local-uniform-convergence hypothesis nor the continuous-dominant hypothesis of Theorem~\ref{thm:dct} holds. Pointwise convergence away from a single point is not enough for the stated theorem.
\end{solution}

\begin{solution}{ex:densityunique}
\idea The fundamental theorem recovers a continuous integrand from its integrals.
\writeup Let $H(x)=\int_0^x(f-g)$. Then $H$ is identically $0$, and $H'=f-g$ by Volume~1, Chapter~8. So $f=g$.
\end{solution}

\begin{solution}{ex:tail}
\idea The Gaussian decay beats any polynomial, and the comparison $t\ge x$ turns the integral into an exact derivative.
\writeup For $t\ge x>0$, $e^{-t^2/2}\le (t/x)e^{-t^2/2}$. Integrate both sides from $x$ to $\infty$. The right-hand side equals $e^{-x^2/2}/x$. Divide by $\sqrt{2\pi}$ to pass from the integrand of $\phi$ to the probability.
\end{solution}

\begin{solution}{ex:cauchy}
\idea The antiderivative is an arctangent. The tails of $|x|f(x)$ are at least a constant times $1/|x|$.
\writeup $\int \mathrm dx/(1+x^2)=\arctan x$, so $\int_{-\infty}^{\infty}\mathrm dx/(\pi(1+x^2))=(1/\pi)(\pi/2-(-\pi/2))=1$. For $|x|\ge1$, $|x|/(1+x^2)\ge |x|/(x^2+x^2)=1/(2|x|)$. The integral of $1/|x|$ from $1$ to $\infty$ diverges, so $\int |x|f=\infty$.
\end{solution}

\begin{solution}{ex:gamma}
\idea Each integration by parts lowers the power by one and produces a factor $n$.
\writeup $I_1=\int_0^{\infty}e^{-x}\,\dd x=1$. For $n\ge1$, integrate $I_{n+1}=\int_0^{\infty}x^n e^{-x}\,\dd x$ by parts: $u=x^n$, $\mathrm dv=e^{-x}\,\dd x$, so $\mathrm du=nx^{n-1}\,\dd x$ and $v=-e^{-x}$. The boundary term $[-x^n e^{-x}]_0^{\infty}$ is $0$, because $x^n e^{-x}\to0$ at infinity by comparing with the series for $e^{x}$. Thus $I_{n+1}=n I_n$, and $I_n=(n-1)!$.
\end{solution}

\begin{solution}{ex:order}
\idea One sample sits in a small bin at $t$, $k-1$ sit below, and the rest sit above. The mean is a ratio of factorials, from one integration by parts at a time.
\writeup The probability that $U_{(k)}$ falls in $(t,t+\Delta t)$ equals
\[
 n\binom{n-1}{k-1}\,(\Delta t)\,t^{k-1}(1-t)^{n-k}
\]
plus an error $o(\Delta t)$ as $\Delta t\to0$. Dividing by $\Delta t$ gives the density on $[0,1]$.

Set $I(k,m)=\int_0^1 t^k(1-t)^m\,\dd t$ for integers $k,m\ge0$. For $k\ge1$, integrate by parts with $u=t^k$ and $\mathrm dv=(1-t)^m\,\dd t$. Then $\mathrm du=kt^{k-1}\,\dd t$ and $v=-(1-t)^{m+1}/(m+1)$. The boundary term vanishes at both ends, and
\[
 I(k,m)=\frac{k}{m+1}I(k-1,m+1).
\]
Iterate until the power of $t$ is gone:
\[
 I(k,m)=\frac{k!}{(m+1)\cdots(m+k)}I(0,m+k)=\frac{k!\,m!}{(m+k+1)!},
\]
because $I(0,r)=\int_0^1(1-t)^r\,\dd t=1/(r+1)$. Therefore
\[
 \int_0^1 t^{k}(1-t)^{n-k}\,\dd t=I(k,n-k)=\frac{k!(n-k)!}{(n+1)!}.
\]
The mean is
\[
 n\binom{n-1}{k-1}\cdot\frac{k!(n-k)!}{(n+1)!}
 =n\cdot\frac{(n-1)!}{(k-1)!(n-k)!}\cdot\frac{k!(n-k)!}{(n+1)!}
 =\frac{k}{n+1}.
\]
\end{solution}

\begin{solution}{ex:cantor}
\idea The series for $X$ is a ternary expansion that never uses the digit $1$. The removed middle thirds are exactly the numbers that are forced to use the digit $1$.
\writeup (a) Stage $0$ is $[0,1]$, length $1$. Each stage removes the open middle third of every remaining closed interval, and leaves two closed intervals of one third the length. After $n$ stages, $2^n$ intervals of length $3^{-n}$ remain, with total length $(2/3)^n\to0$. The set $C$ is contained in the remaining set at every stage.

(b) The number $X$ has ternary expansion $0.(2B_1)(2B_2)(2B_3)\ldots$, because $2B_k\in\{0,2\}$. So $X\in C$ for every outcome. A single value $x$ specifies every ternary digit, except on the countable set of endpoints, which have two expansions. Either way only one sequence of bits, or two, realises $x$. A single infinite bit string has probability $0$, being contained in cylinders of length $2^{-N}$ for every $N$.

(c) We can check continuity without assuming an unproved general continuity theorem for probabilities. Fix $x$ and $\varepsilon>0$. At stage $n$ of the construction, each retained interval carries probability $2^{-n}$, and at most two such intervals can meet a sufficiently short interval around $x$ (their distinct interiors have a positive gap unless they share an endpoint). Choose $n$ with $2^{1-n}<\varepsilon$, then choose the neighbourhood small enough to meet at most two stage-$n$ intervals. Its probability is at most $2^{1-n}<\varepsilon$, proving continuity of $P(X\le x)$. Suppose a continuous $f$ were a density. On each removed middle-third interval the probability is $0$, so $\int f=0$ there, and continuity forces $f=0$ on that interval. The removed intervals are dense in $[0,1]$: every subinterval of $[0,1]$ contains a removed middle third at a late enough stage. A continuous function that vanishes on a dense subset of $[0,1]$ vanishes throughout that interval. The density must also vanish outside $[0,1]$, because each interval wholly outside has probability zero. It would therefore integrate to $0$, contradicting the required total mass $1$.
\end{solution}

\begin{solution}{ex:usquare}
\idea Squaring folds nothing on $(0,1)$, because the square root is the inverse there.
\writeup For $y\in(0,1)$, $P(Y\le y)=P(U\le\sqrt y)=\sqrt y$, since $U$ is uniform on $(0,1)$. The distribution function is $0$ for $y\le0$ and $1$ for $y\ge1$. Differentiating on $(0,1)$ gives the density $1/(2\sqrt y)$. Then
\[
 \mathbb E[Y]=\int_0^1 y\cdot\frac{1}{2\sqrt y}\,\dd y=\int_0^1\frac12 y^{1/2}\,\dd y=\frac13.
\]
The same number is $\mathbb E[U^2]=\int_0^1 u^2\,\dd u=1/3$, which is the definition of the expectation for this bounded continuous function of $U$.
\end{solution}

\begin{solution}{ex:radius}
\idea Box--Muller sets $U=e^{-R^2/2}$, and $U$ is uniform, so $R^2$ is a transformation of a uniform.
\writeup $P(R^2>t)=P(-2\ln U>t)=P(U<e^{-t/2})=e^{-t/2}$ for $t>0$, because $U$ is uniform on $(0,1)$ and $-2\ln u$ is decreasing in $u$. The survival function $e^{-t/2}$ is the exponential law of rate $1/2$, whose density is $\tfrac12 e^{-t/2}$.
\end{solution}

\section*{Sources and further reading}

\begin{enumerate}
\item G.~Vitali, Sul problema della misura dei gruppi di punti di una retta, Bologna, 1905. The obstruction to a length for every subset, under countable additivity and translation invariance.
\item Volume~1, Chapter~8, for the Riemann integral, the fundamental theorem, and $\int_0^{\infty}e^{-u^2}\,\dd u=\sqrt{\pi}/2$.
\item G.~E.~P.~Box and M.~E.~Muller, A note on the generation of random normal deviates, \emph{Ann.\ Math.\ Statist.} \textbf{29} (1958), 610--611. Theorem~\ref{thm:boxmuller}.
\end{enumerate}
